% ---------------- บทคัดย่อภาษาไทย ----------------
\chapter*{บทคัดย่อ}
\addcontentsline{toc}{chapter}{บทคัดย่อ}

กระบวนการค้นหาและพัฒนายาใหม่เผชิญข้อจำกัดด้านระยะเวลาและต้นทุนมหาศาล การทดสอบฤทธิ์ทางชีวภาพของสารเคมีในห้องปฏิบัติการจริงกับโปรตีนเป้าหมาย เช่น เอนไซม์ Main Protease (\Mpro\ หรือ \CLpro) ของเชื้อไวรัส SARS-CoV-2 มีขั้นตอนซับซ้อนและมีค่าใช้จ่ายสูง โครงงานนี้จึงพัฒนาแบบจำลองการเรียนรู้ของเครื่องเพื่อทำนายค่าความสามารถในการยับยั้งเอนไซม์ในรูป \pICfifty\ แบบ In Silico

คณะผู้จัดทำรวบรวมข้อมูลจากฐานข้อมูล ChEMBL เป้าหมาย \texttt{CHEMBL4523582} จำนวน 4{,}867 รายการ กรองเฉพาะ assay ที่วัด \CLpro\ เหลือ 2{,}889 รายการ และผ่านกระบวนการทำความสะอาด 13 ขั้นตอน ได้ชุดข้อมูลมาตรฐาน 2{,}198 โมเลกุลที่ไม่ซ้ำกัน สกัดคุณลักษณะแบบผสม (Hybrid Features) จำนวน 1{,}028 มิติ ประกอบด้วย Morgan Fingerprints (ECFP4 รัศมี 2 ขนาด 1{,}024 บิต) และ Lipinski Descriptors 4 ตัว จากนั้นเปรียบเทียบแบบจำลอง 4 ชนิด ได้แก่ Dummy Regressor, Random Forest, XGBoost และ Support Vector Regressor ผ่าน Scikit-learn Pipeline ที่ป้องกันการรั่วไหลของข้อมูล พร้อมทำ 5-Fold Cross-Validation และปรับจูนไฮเปอร์พารามิเตอร์

ผลการทดลองพบว่า XGBoost ให้ประสิทธิภาพสูงสุดบนชุดทดสอบ โดยมีค่า \Rsq\ = 0.7034, RMSE = 0.5487, MAE = 0.3918 และ Spearman $\rho$ = 0.8415 ซึ่งเหนือกว่าแบบจำลองพื้นฐานอย่างชัดเจน (\Rsq\ $\approx 0.0000$) การทำนายร้อยละ 93.6 มีความคลาดเคลื่อนไม่เกิน 1 log unit เมื่อทดสอบด้วยการแบ่งข้อมูลตามโครงสร้างแกนกลาง (Scaffold Split) ค่า \Rsq\ ลดลงเหลือ 0.5127 ซึ่งสะท้อนความสามารถทำนายที่แท้จริงเมื่อพบสารแกนโครงสร้างใหม่

การวิเคราะห์ด้วยเทคนิค SHAP พบว่าน้ำหนักโมเลกุลเป็นคุณลักษณะสำคัญที่สุด (Mean$|$SHAP$|$ = 0.1165) และสามารถระบุโครงสร้างย่อย 2 มิติที่สัมพันธ์กับการเพิ่มและลดฤทธิ์ยับยั้งได้ 6 บิต โดยบิตที่ 301 พบในสารกลุ่ม active ร้อยละ 50.9 แต่พบในกลุ่ม inactive เพียงร้อยละ 13.1

ข้อค้นพบที่สำคัญคือ การทดสอบข้ามชุดข้อมูลกับสาร COVID-19 จากแหล่งอิสระ 90 ตัว ให้ค่า \Rsq\ = $-1.72$ ซึ่งแย่กว่าการทำนายด้วยค่าเฉลี่ย สาเหตุมาจากร้อยละ 84.4 ของสารกลุ่มดังกล่าวอยู่นอกขอบเขตการใช้งานของแบบจำลอง (Applicability Domain) ผลนี้ชี้ว่าแบบจำลอง QSAR ที่มีค่า \Rsq\ สูงบนชุดทดสอบภายในไม่ได้รับประกันความสามารถทำนายสารนอกปริภูมิเคมีเดิม

\vspace{0.8em}
\noindent\textbf{คำสำคัญ:} การเรียนรู้ของเครื่อง, QSAR, SARS-CoV-2 Main Protease, \pICfifty, Morgan Fingerprints, SHAP, Applicability Domain
